% !TEX program = pdflatex
% Presentation superviseur -- Week 7 -- 2026-05-26
% Étape 1 : embedding articles-first + résultats comparés vs baseline 5 mai

\documentclass[aspectratio=169,10pt]{beamer}

\usepackage[utf8]{inputenc}
\usepackage[T1]{fontenc}
\usepackage[french]{babel}
\usepackage{lmodern}
\usepackage{booktabs}
\usepackage{array}
\usepackage{tikz}
\usepackage[table]{xcolor}
\usepackage{graphicx}
\usetikzlibrary{positioning, arrows.meta, shapes.geometric, fit, backgrounds}

\usetheme{Madrid}
\usecolortheme{seahorse}
\setbeamertemplate{navigation symbols}{}
\setbeamertemplate{footline}[frame number]
\setbeamerfont{footnote}{size=\tiny}

\definecolor{accent}{RGB}{41, 92, 155}
\definecolor{accentsoft}{RGB}{230, 238, 247}
\definecolor{warn}{RGB}{200, 80, 40}
\definecolor{ok}{RGB}{30, 130, 70}
\definecolor{ppcol}{RGB}{47, 111, 79}
\definecolor{penalcol}{RGB}{124, 45, 45}

\newcommand{\keybox}[1]{%
  \begin{center}
  \colorbox{accentsoft}{\parbox{0.92\linewidth}{\centering #1}}
  \end{center}}
\newcommand{\warnbox}[1]{%
  \begin{center}
  \colorbox{red!10}{\parbox{0.92\linewidth}{\centering #1}}
  \end{center}}
\newcommand{\okbox}[1]{%
  \begin{center}
  \colorbox{green!10}{\parbox{0.92\linewidth}{\centering #1}}
  \end{center}}

\title[KG juridique FR --- Week 7]{Construction d'un Knowledge Graph juridique français}
\subtitle{Point d'avancement --- Week 7 : Étape 1, embedding articles-first}
\author{Matthieu Kaeppelin}
\institute{Stage FE Recherche}
\date{26 mai 2026}

\begin{document}

% ============================================================
\begin{frame}
  \titlepage
\end{frame}

% ============================================================
\begin{frame}{Plan de la séance}
  \tableofcontents[hideallsubsections]
\end{frame}

% ============================================================
\section{Rappel : de Week 4 à Week 7}

\begin{frame}{Où en étions-nous à Week 4 (5 mai) ?}
  \small
  \textbf{Constat principal de Week 4} :

  \keybox{La baseline B2 (embedding \texttt{multilingual-e5-base} sur les arrêts) ne
  ressort \emph{aucune} référence pertinente avant le top-1000. Le bottleneck est
  l'\textbf{encodeur}, pas le retrieval.}

  \vspace{0.5em}
  \textbf{Chiffres clés du 5 mai} :
  \begin{itemize}\setlength\itemsep{0.2em}
    \item Rang médian des arrêts gold : \textcolor{warn}{\textbf{76 782 / 118 112}}
    \item Recall JP top-10 (toutes stratégies) : \textcolor{warn}{\textbf{0 / 16}}
    \item Couverture gold dans le corpus : \textcolor{ok}{100 \%} (donc problème de retrieval, pas de données)
  \end{itemize}

  \vspace{0.5em}
  \textbf{Roadmap actée avec Johnny} (point S4 §7) :
  \begin{itemize}\setlength\itemsep{0.2em}
    \item \textbf{Étape 1} : embedding pur sans graphe (articles + JP), mesurer recall@K
    \item Étape 2 : combler les JP sans summary (LLM ou laplacien)
    \item Étape 3 : benchmark étendu via doctrine + GNN link prediction
  \end{itemize}
\end{frame}

\begin{frame}{Roadmap actualisée}
  \footnotesize
  \renewcommand{\arraystretch}{1.15}
  \begin{tabular}{@{}c p{4.4cm} p{5.6cm} l@{}}
    \toprule
    \textbf{\#} & \textbf{Palier} & \textbf{Question scientifique} & \textbf{État} \\
    \midrule
    P1 & Extraction d'articles       & Sait-on identifier les articles cités ?              & \textcolor{ok}{\textbf{Validé S3}} \\
    P2 & Robustesse extraction       & LLMs vs regex symbolique ?                            & \textcolor{ok}{\textbf{Validé S3}} \\
    P3 & Construction graphe complet & Pipeline tient-il sur 1,13 M arrêts ?                & \textcolor{ok}{\textbf{Validé S3}} \\
    P4 & Analyse macro du graphe     & Quelle structure émerge ?                             & \textcolor{ok}{\textbf{Validé S3}} \\
    \rowcolor{accentsoft}
    P5 & Embedding articles (LEGI)   & BGE-M3 sur articles bat-il e5-base sur JP ?           & \textcolor{accent}{\textbf{Validé S7}} \\
    \rowcolor{accentsoft}
    P6 & Évaluation recall@K         & Quel régime de performance atteint-on ?               & \textcolor{accent}{\textbf{Validé S7}} \\
    P7 & Complétion JP (Étape 2)     & LLM ou laplacien complète-t-il le signal ?            & \textcolor{warn}{\textbf{Next}} \\
    P8 & GNN link prediction         & Le graphe enrichi bat-il le retrieval pur ?           & Étape 3 \\
    \bottomrule
  \end{tabular}

  \vspace{0.8em}
  \keybox{\textbf{Cette session} : présenter les résultats P5+P6 et leur diagnostic.}
\end{frame}

% ============================================================
\section{Étape 1 : design et pipeline}

\begin{frame}{Le pivot stratégique de l'Étape 1}
  \small
  \textbf{Phase D du journal 5 mai} : \emph{« et si on embeddait les \textbf{articles},
  pas les arrêts ? »}

  \vspace{0.5em}
  \textbf{Justification a priori} :
  \begin{itemize}\setlength\itemsep{0.2em}
    \item Une question CRFPA cible un \emph{point de droit codifié} (article), pas un narratif d'arrêt
    \item Articles = textes courts, denses, structurés sémantiquement
    \item Arrêts = textes longs, narratifs, mêlent faits + procédure + motifs + dispositif (bruit)
    \item Mapping question $\to$ article \emph{plus direct} que question $\to$ arrêt
  \end{itemize}

  \vspace{0.5em}
  \textbf{Objection à lever} : on n'a pas le texte des articles -- juste les
  \texttt{pair\_key} (\texttt{code\_penal:222-23}) extraits par regex V5.

  \vspace{0.5em}
  \okbox{\textbf{Solution} : ingestion du dump LEGI ouvert (DILA) via \texttt{legi.py}.
  Résolution \texttt{pair\_key} $\to$ texte intégral.}
\end{frame}

\begin{frame}{Pipeline en 4 étages}
  \begin{center}
  \begin{tikzpicture}[node distance=0.8cm and 1.6cm,
    box/.style={draw, rectangle, rounded corners=4pt, minimum width=2.8cm, minimum height=1cm, align=center, font=\scriptsize},
    accentbox/.style={box, fill=accentsoft, draw=accent},
    arrow/.style={-{Stealth[length=2.5mm]}, thick}]

    \node[accentbox] (s1) {\textbf{1. Diagnostic}\\Token-stats\\(p100 articles + JP)};
    \node[accentbox, right=of s1] (s2) {\textbf{2. Ingestion LEGI}\\Dump DILA $\to$ SQLite\\\texttt{pair\_key} $\to$ texte};
    \node[accentbox, right=of s2] (s3) {\textbf{3. Embedding}\\BGE-M3 1024d\\(articles)};
    \node[accentbox, below=of s3] (s4) {\textbf{4. Évaluation}\\recall@K\\question $\to$ articles};
    \node[box, dashed, draw=warn, below=of s1, font=\tiny] (linkage) {Linkage \\\texttt{pair\_key} $\leftrightarrow$ graphe};
    \node[box, dashed, draw=warn, below=of s2, font=\tiny] (jpgraph) {Back-edge graphe\\top-K art $\to$ JP citantes};

    \draw[arrow] (s1) -- (s2);
    \draw[arrow] (s2) -- (s3);
    \draw[arrow] (s3) -- (s4);
    \draw[arrow] (s2) -- (linkage);
    \draw[arrow] (linkage) -- (s4);
    \draw[arrow] (s4) -- (jpgraph);
  \end{tikzpicture}
  \end{center}

  \vspace{0.3em}
  \footnotesize
  \textbf{Invariant central} : la clé \texttt{pair\_key} relie nœud du graphe, texte LEGI
  et ligne d'embedding. Alignement positionnel \emph{lossless} $\to$ l'Étape 2 (laplacien
  sur le graphe) récupère les embeddings sans jointure a posteriori.
\end{frame}

\begin{frame}{Paramètres expérimentaux}
  \footnotesize
  \begin{tabular}{@{}p{4cm} p{8cm}@{}}
    \toprule
    \textbf{Composant} & \textbf{Configuration} \\
    \midrule
    Modèle d'embedding & \texttt{BAAI/bge-m3} (1024d, 8k tok théoriques, multilingue) \\
    Contexte effectif & 2 048 tok (limite pratique MPS, chunk+meanpool au-delà) \\
    Source des articles & Dump DILA LEGI, snapshot 13 juillet 2025 (1,73 M articles tous codes) \\
    Pool candidat & 3 183 articles pénaux résolus (subset des 8 085 nœuds du graphe) \\
    Couverture LEGI & 39 \% globale, \textbf{100 \% sur le gold} obligatoire (50/50) \\
    Hardware & Mac MPS, batch 8, $L_2$ normalisation \\
    Questions & 8 CRFPA pénales (3 droit pénal + 5 procédure pénale) \\
    Métrique principale & recall@K, $K^* = $ plus petit $K$ tel que recall $\geq 0,5$ \\
    \bottomrule
  \end{tabular}

  \vspace{0.6em}
  \keybox{Le pool candidat (3 183 vs 118 112) est \emph{37$\times$ plus petit} que la
  baseline B2 -- l'éval mécaniquement plus indulgente.}
\end{frame}

% ============================================================
\section{Résultats : critère strict article-first}

\begin{frame}{Métrique adoptée pour pénaliser la sur-extraction}
  \small
  \textbf{Problème de la métrique recall@1000} : on peut « gagner » en remontant tout le
  pool, donc trop laxiste pour mesurer une exploitation réelle.

  \vspace{0.5em}
  \textbf{Métriques bornées top-K adoptées} :

  \vspace{0.3em}
  \begin{tabular}{@{}p{3cm} p{4.5cm} p{4cm}@{}}
    \toprule
    & \textbf{Articles} & \textbf{JP via graphe} \\
    \midrule
    Critère \textbf{dur} & recall@10 $\geq$ 0,5 & recall@5 $\geq$ 0,5 \\
    Critère \textbf{facile} & recall@20 $\geq$ 0,5 & recall@10 $\geq$ 0,5 \\
    \bottomrule
  \end{tabular}

  \vspace{0.5em}
  \keybox{Question passe le critère $\Longleftrightarrow$ les \emph{deux} conditions
  simultanément. Équivalent d'un budget de 10/5 résultats affichables.}

  \vspace{0.5em}
  \textbf{Bornage implicite de la précision} :
  pour 1 article gold, recall@10 = 1 $\to$ précision max 10 \%.
  Pour 7 articles gold, recall@10 = 1 $\to$ précision max 70 \%.
  Donc la métrique mesure aussi le \emph{non-noyage du gold dans le bruit}.
\end{frame}

\begin{frame}{Résultats par question (critère dur)}
  \footnotesize
  \renewcommand{\arraystretch}{1.2}
  \begin{tabular}{@{}l c c c c c c c@{}}
    \toprule
    \textbf{Question} & \textbf{Branche} & \textbf{gold o} & \textbf{gold JP} &
    \textbf{r@10 art} & \textbf{r@5 JP} & \textbf{Passe ?} & \textbf{Diagnostic} \\
    \midrule
    PP-Q2 & \textcolor{ppcol}{PP} & 1 & 2 & \textbf{1.00} & \textbf{0.50} & \textcolor{ok}{\textbf{Oui}} & idéal \\
    PP-Q3 & \textcolor{ppcol}{PP} & 2 & 1 & \textbf{1.00} & \textbf{1.00} & \textcolor{ok}{\textbf{Oui}} & idéal \\
    PP-Q1 & \textcolor{ppcol}{PP} & 2 & 3 & \textbf{0.50} & \textbf{0.67} & \textcolor{ok}{\textbf{Oui}} & 1 art/2 trouvé \\
    PP-Q4 & \textcolor{ppcol}{PP} & 1 & 1 & \textbf{1.00} & 0.00 & \textcolor{warn}{\textbf{Non}} & art OK, graphe JP fail \\
    PP-Q5 & \textcolor{ppcol}{PP} & 2 & 0 & \textbf{0.50} & n/a & \textcolor{warn}{$\circ$} & gold JP absent \\
    PENAL-Q1 & \textcolor{penalcol}{PEN} & 7 & 8 & 0.14 & 0.12 & \textcolor{warn}{\textbf{Non}} & contamination thématique \\
    PENAL-Q3 & \textcolor{penalcol}{PEN} & 2 & 2 & 0.00 & 0.00 & \textcolor{warn}{\textbf{Non}} & principe abstrait raté \\
    PENAL-Q2 & \textcolor{penalcol}{PEN} & 2 & 2 & 0.00 & 0.00 & \textcolor{warn}{\textbf{Non}} & compositionnel pur \\
    \bottomrule
  \end{tabular}

  \vspace{0.8em}
  \okbox{\textbf{3 / 7 questions évaluables passent le critère dur.}
  Baseline 5 mai : \textbf{0 / 7}. PP-Q5 est exclue (gold JP absent du corpus).}
\end{frame}

\begin{frame}{Comparaison vs baseline B2 (5 mai)}
  \footnotesize
  \renewcommand{\arraystretch}{1.25}
  \begin{tabular}{@{}p{4cm} p{3.6cm} p{3.6cm} l@{}}
    \toprule
    \textbf{Métrique} & \textbf{B2 -- 5 mai} & \textbf{Étape 1 -- 26 mai} & $\Delta$ \\
    \midrule
    Modèle & e5-base 512 tok & \textbf{BGE-M3 2k tok} & + 4$\times$ contexte \\
    Entrée encodée & texte JP & \textbf{texte article LEGI} & inversion \\
    Pool candidat & 118 112 JP & 3 183 articles & $\times$ 37 plus petit \\
    \midrule
    mean recall@10 art (oblig) & 0,158 & \textbf{0,518} & $\times$ 3,3 \\
    mean recall@5 JP & \textcolor{warn}{\textbf{0,000}} & \textbf{0,286} & $\infty$ \\
    mean recall@1000 JP & 0,12 (12 \%) & 0,94 & $\times$ 7,8 \\
    Rang médian gold JP & 76 782 / 118 112 & $\leq$ 100 (majorité) & -- \\
    \midrule
    \textbf{Questions passant critère dur} & \textcolor{warn}{\textbf{0 / 7}} & \textbf{3 / 7} & saut qualitatif \\
    \bottomrule
  \end{tabular}

  \vspace{0.5em}
  \keybox{Pas un gain quantitatif (\og un peu mieux \fg{}) mais un \textbf{changement de
  régime} : de \og rien ne sort \fg{} à \og 3 questions utilisables sur 7 \fg{}.}
\end{frame}

% ============================================================
\section{Diagnostic par question}

\begin{frame}{Le pattern central : PP $\gg$ PENAL}
  \footnotesize
  \begin{tabular}{@{}l c c c@{}}
    \toprule
    & \textbf{$K^*$ médian articles} & \textbf{questions critère dur} & \textbf{interprétation} \\
    \midrule
    \textcolor{ppcol}{\textbf{Procédure pénale}} (5 questions) & \textbf{3} & 3/4 évaluables & \textcolor{ok}{lexical} \\
    \textcolor{penalcol}{\textbf{Droit pénal de fond}} (3 questions) & \textbf{200} & 0/3 & \textcolor{warn}{compositionnel} \\
    \bottomrule
  \end{tabular}

  \vspace{0.8em}
  \textbf{Hypothèse explicative} :

  \begin{itemize}\setlength\itemsep{0.3em}
    \item \textbf{Procédure pénale} = requêtes lexicales sur articles bien identifiés.
    \og garde à vue \fg{} $\to$ CPP:63. \og IPC + nullité \fg{} $\to$ CPP:173-1. Le
    vocabulaire de la question matche celui de l'article -- \emph{terrain idéal pour
    l'embedding}.
    \item \textbf{Droit pénal de fond} = cas pratiques compositionnels. Un récit
    factuel (\og dirigeant, médicament, retard d'alerte \fg{}) qui exige une
    \emph{qualification juridique} (\og blessures involontaires + manquement de
    prudence \fg{}) mobilisant des articles-principes abstraits (\texttt{121-3},
    \texttt{121-1}) jamais nommés dans la question.
  \end{itemize}

  \vspace{0.3em}
  \warnbox{\textbf{L'embedding lexico-sémantique ne raisonne pas.} Il matche
  thématiquement, mais ne fait pas d'inférence de qualification.}
\end{frame}

\begin{frame}{Trois patterns d'échec, trois fixes différents}
  \footnotesize
  \renewcommand{\arraystretch}{1.4}
  \begin{tabular}{@{}p{2.6cm} p{2.4cm} p{4.2cm} p{4cm}@{}}
    \toprule
    \textbf{Pattern} & \textbf{Questions} & \textbf{Diagnostic} & \textbf{Fix plausible} \\
    \midrule
    \textbf{A.} Compositionnel pur & PENAL-Q2 (partiel Q1) & Aucun mot-clé gold dans la question
    & Réécriture LLM (\og dirigeant non-alerté \fg{} $\to$ \og imprudence \fg{}) \\
    \textbf{B.} Article-principe raté & PENAL-Q3 (\texttt{121-1}) & Article trop générique,
    non thématique & Multi-vecteur (un par qualification attendue) \\
    \textbf{C.} Back-edge graphe & PP-Q4 (article OK, JP fail) & Article gold trouvé mais
    pourvoi gold ne cite pas cet article & Élargir top-K par articles voisins (2-hop) \\
    \bottomrule
  \end{tabular}

  \vspace{0.6em}
  \keybox{Les 3 patterns appellent des \textbf{fixes différents}. L'Étape 2 (complétion JP) ne résoudra
  que A et indirectement C. B nécessite une stratégie distincte (multi-vecteur ou liste
  d'articles-principes).}
\end{frame}

\begin{frame}{Anatomie d'un échec : PENAL-Q2}
  \footnotesize
  \textbf{Question} : \emph{Thierry, dirigeant du laboratoire 'Medica SA' fabricant le
  médicament prescrit à Romuald, a été informé des effets indésirables dès fin 2023
  mais n'a alerté les médecins et patients qu'en...}

  \textbf{Gold OBLIG} : \texttt{code\_penal:222-19} (blessures involontaires),
  \texttt{code\_penal:121-3} (intention / imprudence).

  \vspace{0.4em}
  \textbf{Rangs des gold dans le classement} :
  \begin{itemize}\setlength\itemsep{0em}
    \item \texttt{code\_penal:222-19} -- rang \textcolor{warn}{\textbf{836}} / 3 183
    \item \texttt{code\_penal:121-3} -- rang \textcolor{warn}{\textbf{971}} / 3 183
  \end{itemize}

  \vspace{0.3em}
  \textbf{Top-5 retourné par BGE-M3} :
  \begin{enumerate}\setlength\itemsep{0em}
    \item \texttt{CPP:706-47-1} -- suivi des condamnés pour infractions sexuelles
    \item \texttt{CPP:D47-27} -- troubles mentaux / détention
    \item \texttt{CPP:706-136} -- chambre instruction, décision
    \item \texttt{CP:222-33} -- harcèlement sexuel
    \item \texttt{CP:222-28} -- agression sexuelle
  \end{enumerate}

  \vspace{0.3em}
  \warnbox{Le mot-clé décisif (\og blessures involontaires \fg{} ou \og élément moral \fg{})
  n'est nulle part dans la question. Il est dans la \emph{qualification juridique} qu'un
  correcteur applique sur le cas pratique. \textbf{L'embedding ne fait pas ce raisonnement}.}
\end{frame}

\begin{frame}{Anatomie d'une réussite : PP-Q3}
  \footnotesize
  \textbf{Question} : \emph{À la suite de leur interrogatoire de première comparution
  le 3 mars 2025, le juge d'instruction a mis en examen les trois suspects y compris
  ...}

  \textbf{Gold OBLIG} : \texttt{CPP:116-1}, \texttt{CPP:116}.

  \vspace{0.4em}
  \textbf{Top-5 retourné par BGE-M3} :
  \begin{enumerate}\setlength\itemsep{0em}
    \item \texttt{CPP:116-1} -- interrogatoires criminels (\textcolor{ok}{\textbf{OBLIG}})
    \item \texttt{CPP:80-1} -- mise en examen
    \item \texttt{CPP:116} -- interrogatoire IPC (\textcolor{ok}{\textbf{OBLIG}})
    \item \texttt{CPP:176} -- examen du juge d'instruction
    \item \texttt{CPP:113-8} -- mise en examen, statut témoin assisté
  \end{enumerate}

  \vspace{0.3em}
  \okbox{Les 2 articles gold sont en rangs 1 et 3. Recall@10 articles = 1.00.
  Recall@5 JP via graphe = 1.00. Cas idéal : question lexicale, articles codifiés,
  citations JP nettes.}
\end{frame}

% ============================================================
\section{Limitations et travail en cours}

\begin{frame}{Limitations chiffrées}
  \footnotesize
  \begin{tabular}{@{}p{3.8cm} p{2cm} p{7cm}@{}}
    \toprule
    \textbf{Limitation} & \textbf{Chiffre} & \textbf{Impact} \\
    \midrule
    Couverture LEGI globale & 39 \% & Articles pré-1994 renumérotés. Gold à 100 \%, éval valide. \\
    Échantillon questions & 8 & Trop peu pour stat fiable sur PP vs PENAL. Étape 3. \\
    Politique troncature & 2 048 tok & 5 articles (0,16 \%) chunkés. Décision pragmatique OOM MPS. \\
    JP gold extractibles & CC-only & \texttt{ref:null} dans rubrics. 1/8 questions sans gold JP. \\
    Gold orphelin graphe & 1/51 & \texttt{code\_penal:222-26-2} non cité par aucune JP. \\
    \bottomrule
  \end{tabular}

  \vspace{0.6em}
  \keybox{\textbf{Aucune limitation n'invalide la conclusion principale} (changement de
  régime). Elles définissent le \emph{plafond} de l'approche, pas son existence.}
\end{frame}

\begin{frame}{En cours : élargissement au full corpus}
  \small
  \textbf{Question naturelle} : et si on embeddait les \textbf{87 821} articles du graphe
  (tous codes, pas juste les 4 pénaux) ?

  \vspace{0.5em}
  \textbf{Pourquoi} :
  \begin{itemize}\setlength\itemsep{0.2em}
    \item Tester si les questions pénales \emph{tirent} des articles non-pénaux pertinents (code civil pour les questions de fond, par exemple).
    \item Mesurer la robustesse du classement quand le pool est $\times$ 27 plus grand.
    \item Vérifier si le critère dur tient sur 87 k candidats.
  \end{itemize}

  \vspace{0.5em}
  \textbf{Estimation} : embedding $\approx$ 2-3 h sur Mac MPS, à $\sim$ 12 MB pour 3 183 articles $\to$ $\sim$ 360 MB pour 87 k.

  \vspace{0.4em}
  \keybox{\textbf{Lancé en background pendant cette présentation.} Résultats disponibles
  pour la prochaine session.}
\end{frame}

% ============================================================
\section{Prochaines étapes}

\begin{frame}{Étape 2 : complétion JP via le graphe}
  \small
  \textbf{Objectif} : produire des embeddings JP exploitables pour les questions
  compositionnelles (patterns A et C).

  \vspace{0.4em}
  \textbf{Deux méthodes en parallèle} (cf. note du 5 mai sur la complétion par
  optimisation laplacienne) :
  \begin{itemize}\setlength\itemsep{0.3em}
    \item \textbf{Méthode (A)} -- batch LLM : générer un summary synthétique pour
    chaque JP sans summary, embedder ce summary.
    \item \textbf{Méthode (B)} -- complétion laplacienne :
    $H^\star = \arg\min_H \lVert \widetilde{H} - H \odot J \rVert_F^2 + \lambda \cdot \mathrm{tr}(H^\top L H)$,
    propage l'info des JP avec summary vers les JP sans, via le graphe.
  \end{itemize}

  \vspace{0.4em}
  \textbf{Diagnostic préalable obligatoire} : énergie de Dirichlet
  $\rho(\widetilde{H}_o) = \mathrm{tr}(\widetilde{H}_o^\top L_{\mathrm{sym}}^{(o)} \widetilde{H}_o) / \lVert \widetilde{H}_o \rVert_F^2$
  + test de permutation. \textbf{Si le graphe ne porte pas de signal, repli sur (A).}

  \vspace{0.4em}
  \keybox{\textbf{Critère d'acceptation Étape 2} : passer de \textbf{3/7} à
  $\geq$ \textbf{5/7} questions sur le critère dur, en particulier sur PENAL-Q1 et
  PENAL-Q2.}
\end{frame}

\begin{frame}{Étape 3 : benchmark étendu via doctrine}
  \small
  \textbf{Problème} : 8 questions = échantillon trop faible pour distinguer
  fiablement le signal du bruit, surtout sur le pattern PENAL.

  \vspace{0.4em}
  \textbf{Pipeline proposé (cf. roadmap Johnny)} :
  \begin{enumerate}\setlength\itemsep{0.2em}
    \item Collecte de doctrine pénale (manuels, revues juridiques)
    \item Extraction de questions par LLM avec validation croisée 2 LLM locaux \emph{(déjà en cours côté
    pipeline \texttt{doctrine\_qgen}, S5--S6)}
    \item Cible : 100--500 questions étiquetées avec articles + JP de référence
    \item \textbf{Stratification} par type (lexical vs compositionnel) -- alignée sur le
    pattern PP/PENAL identifié à S7
    \item Intégration au KG comme nouveau type de nœud \texttt{Question}
    \item Entraînement GNN (GraphSAGE / R-GCN) en link prediction
  \end{enumerate}

  \vspace{0.4em}
  \keybox{Le \textbf{plancher} étape 1 servira de baseline contre laquelle mesurer le
  GNN. Si le GNN ne bat pas 3/7 sur le critère dur, ce n'est pas un succès.}
\end{frame}

\begin{frame}{Récapitulatif et décisions à arbitrer}
  \small
  \textbf{Acquis de l'Étape 1} :
  \begin{itemize}\setlength\itemsep{0.2em}
    \item Pipeline reproductible, code testé (26/26 tests), 6 commits sur \texttt{etape1-embedding-pur}
    \item 12 MB d'embeddings articles (\texttt{emb\_articles.npy}) -- réutilisable pour les Étapes 2/3
    \item \textbf{3/7 questions} passent le critère dur (recall@10 art $\geq$ 0,5 ET recall@5 JP $\geq$ 0,5)
    \item Diagnostic par question avec 3 patterns d'échec distincts
    \item Documentation à 3 niveaux : pédagogique HTML, technique \texttt{PLAN.md}, synthèse expérimentale
  \end{itemize}

  \vspace{0.4em}
  \textbf{Décisions à arbitrer} :
  \begin{enumerate}\setlength\itemsep{0.2em}
    \item \textbf{Priorité Étape 2 (A) ou (B)} ? La méthode laplacienne est mathématiquement élégante
    mais demande le diagnostic préalable.
    \item \textbf{Multi-vecteur pour les questions compositionnelles} (pattern A/B) -- avant ou après Étape 2 ?
    \item \textbf{Réécriture LLM des questions} comme étape intermédiaire \emph{quasi gratuite} -- à tester en
    parallèle de l'Étape 2 ?
  \end{enumerate}

  \vspace{0.4em}
  \okbox{\textbf{Le plancher est solide. Le travail est sur la trajectoire prévue.}
  Reste à choisir l'angle d'attaque pour faire passer PENAL.}
\end{frame}

\begin{frame}
  \begin{center}
    \vspace{2em}
    {\Large\textbf{Merci.}}
    \vspace{2em}

    \textit{Discussion / questions}

    \vspace{2em}
    \footnotesize
    Branche : \texttt{etape1-embedding-pur} (6 commits)\\
    Synthèse complète : \texttt{06-Analyses/syntheses/Experience-Etape1-Articles-First-2026-05-26.md}\\
    Page résultats HTML : \texttt{01-Projet/presentations/Resultats-Etape1-Articles-First-2026-05-26.html}
  \end{center}
\end{frame}

\end{document}
